\minitoc

% =======================================================================================================
% Introduction 

Dans ce chapitre, nous nous intéressons aux problèmes d'optimisation combinatoire "difficiles"
pour lesquels nous n'avons pas de méthode explicite pour déterminer un solution. 
Nous essayerons de "trier" ces problèmes en classes de complexité et nous verrons que 
certains problèmes sont étrangement similaires les uns des autres.

% =======================================================================================================
% Généralités

\section{Généralités}

Un problème d'optimisation combinatoire est définit par plusieurs éléments. On définit 
ainsi $ \mathcal{X}$ un ensemble de \emph{variables, d"éléments} ou \emph{d'objets}. 
Il représente ainsi l'ensemble des paramètres. La résolution de ce problème vise à 
trouver un \emph{groupement} ou des affectations de $ \mathcal{X}$ dont on $ \mathcal{E}$
l'ensemble des solutions possibles. 

Notre problème peut aussi être définit par un certain nombre de \emph{contraintes} d'où l'existence
d'un ensemble $ \mathcal{S}$ regroupant toutes les solutions \emph{faisables} (i.e respectant les contraintes).

\begin{remark}
    On a donc en général $ \mathcal{S} \subset \mathcal{E}$. $ \mathcal{E}$, énumérer $ \mathcal{S}$ peut 
    être impossible d'où la difficulté du problème. 
\end{remark}

Enfin, nous pouvons aussi définir une fonction de coùt $v : \mathcal{E} \longrightarrow \R$ définissant 
une valeurs ou un coût pour chaque affectation de $ \mathcal{X}$ permettant ainsi d'ordonner les solutions 
et d'en définir un \emph{optimale}.

\begin{definition}[Instance]
    Un \emph{instance} $I$ d'un problème d'optimisation combinatoire est donc un cas particulier 
    de ces problèmes dans un contexte donnée. Elle est caractérisée par un ensemble d'objets/de variables 
    $ \mathcal{X}$ et des paramètres.
\end{definition}

\begin{definition}[Problème d'optimisation combinatoire]
    Ainsi, un \emph{problème d'optimisation combinatoire} noté $\Pi$ est définit par : 
    \begin{itemize}
        \item un ensemble d'instances $ \mathcal{I}$.
        \item pour chaque $ I \in \mathcal{I}$, un ensemble de solutions faisables $ \mathcal{S}(I)$.
        \item une fonction de coût/valeur.
    \end{itemize}
\end{definition}

\begin{remark}
    La différence fondamentale entre une instance et un problème est qu'une instance s'applique dans un 
    contexte donné et a une (ou plusieurs) solution(s) bien spécifique(s).
    Un problème est plutôt la représentation théorique de ce genre de problèmes. 
\end{remark}

Pour un problème donné $ \Pi$ il est intéressant de connaître sa classe de complexité pour savoir comment 
le résoudre et si c'est possible.

Les classes de complexité nous permettront donc de réunir ensemble des problèmes d'optimisation combinatoire 
ayant la même complexité temporelle (dans le pire des cas). 

% =======================================================================================================
% Problèmes de décisions, classes P, NP et NP-complet

\section{Problèmes de décisions, classes P, NP et NP-complet}

\subsection{Complexité temporelle}

\begin{definition}[Problème de décision]
    Un \emph{problème de décision} est un problème dans lequel on souhaite 
    répondre par oui ou par non pour une instance d'un problème donné.

    \medskip 

    On dit qu'une instance d'un problème est \emph{positive} si la réponse au problème 
    de décision est "oui".
\end{definition}

\begin{example}
    Le problème les cliques dans un graphe est un problème de décision. En effet, on souhaite répondre par 
    oui ou non à la question "Un graphe G contient-il un sous-graphe complet d'ordre $k \in \N$ ?
\end{example}

Pour différencier la difficulté de tels problèmes, nous nous intéressons à leur \emph{complexité temporelle}.

\begin{definition}[Complexité temporelle]
    Pour un algorithme donné $A$, on définit sa \emph{complexité temporelle} notée $ Time_A$
    où $ Time_A(x)$ est le nombre d'instructions exécutées par $A$ pour une entrée $x$.
\end{definition}

\begin{proposition}
    On souhaiterai ainsi définir la complexité d'un problème $ \Pi$ par la complexité temporelle 
    de l'algorithme le plus efficace pour résoudre $ \Pi$. Or, il est impossible de dire qu'un algorithme 
    est optimal et ne peut être amélioré. On parlera donc de complexité \emph{asymptotique}.
\end{proposition}

\begin{definition}[Ordre de complexité]
    Soit $f : \R \longrightarrow \R$ une application. Soit $ \Pi$ un problème. On dit que la complexité 
    temporelle de $ \Pi$ est d'ordre $ O(f)$ ssi il existe un algorithme $A$ résulvant $ \Pi$ et 
    $n_0 \in \N, c \in \R$ tels que : 
        \begin{equation}
            \forall x, \quad |x| > n_0 \quad \Longrightarrow \quad Time_A(x) \leqslant c f(|x|)
        \end{equation}
\end{definition}

\subsection{Classe P}

Comme dit précédement, nous souhaitons regrouper des problèmes de même ordre de 
complexité temporelle. Nous définissons donc le concept de classe. 

\begin{definition}[Classes de complexité]
    La classe de complexité $ TIME[f]$ est l'ensemble des problèmes de complexité temporelle $ O(f)$.
    On a donc : 
        $$ TIME[n] \subset TIME[n \log n] \subset TIME[n^2] \subset TIME[e^n] \subset TIME[2^n] $$
\end{definition}

\begin{definition}[Classe P]
    On définit ainsi la classe $P$ comme l'ensemble des classes de problèmes de complexité polynomiale. 
    D'où : 
        \begin{equation}
            \boxed{P = \bigcup_{k \geqslant 0} TIME[n^k] }
        \end{equation}
\end{definition}

La classe $P$ est indépendante du modèle d'ordinateur considéré (machine de Turing ou à RAM). 
Elle est aussi indépendante du langage utilisé pour l'algorithme ainsi que des structures de données 
utilisées pour stocker les variables du problème.

\begin{remark}
    Un problème est donc dit polynomial si il peut être résolut par un algorithme de complexité temporelle 
    polynomiale. Or pour de très grandes valeurs de $k$, cela revient à résoudre un problème de complexité 
    exponentielle. En pratique, nous pourrons donc résoudre des problèmes de complexité $O(n^2), O(n^3)$ ou 
    $O(n^4)$.
\end{remark}

\begin{rappel}[k-SAT]
    Soit $k \in \mathbb{N}$. Un problème de $k$-SAT est un problème de 
    \emph{satisfaisabilité booléenne} portant sur des formules booléennes en forme normale conjonctive (CNF).

    Une formule est composée :
    \begin{itemize}
        \item de \emph{variables booléennes} $x_1, \dots, x_n$ ;
        \item de \emph{littéraux}, qui sont soit une variable soit sa négation : $l \in \{x_i, \lnot x_i\}$ ;
        \item de \emph{clauses}, qui sont des disjonctions de $k$ littéraux :
        \[
            C = l_1 \vee l_2 \vee \dots \vee l_k ;
        \]
        \item et d'une \emph{formule} CNF, qui est une conjonction de clauses :
        \[
            \varphi = \bigwedge_{j = 1}^m C_j .
        \]
    \end{itemize}

    Le problème de $k$-SAT consiste à déterminer s'il existe une assignation de vérité des variables
    qui rende la formule $\varphi$ vraie. Si une telle assignation existe, on dit que la formule est
    \emph{satisfaisable}.
\end{rappel}

\begin{theorem}[SAT]
    2-SAT $ \in P$.
\end{theorem}

\begin{definition}[Réduction polynomiale]
    On dit qu'il existe une \emph{réduction polynomiale} entre deux problèmes $ \Pi_1$ et $ \Pi_2$ si il 
    existe une fonction $f$ calculable en temps polynomial si, $x$ est une instance positive de $ \Pi_1$ 
    ssi $f(x)$ est une instance positive de $ \Pi_2$. On note alors $ \Pi_1 \leqslant \Pi_2$.
\end{definition}

\begin{theorem}[Réduction]
    2-COLORATION $ \leqslant$ 2-SAT.
\end{theorem}


\subsection{Vérifieurs et classe NP}




